\newpage 

\section{Hamiltonian and Lagrangian Mechanics}

\begin{abstract}

\noindent We explore foundational structures of Hamiltonian and Lagrangian mechanics, emphasizing their geometric underpinnings. We begin by defining Hamiltonian systems through a symplectic manifold $(N, \omega)$ representing the phase space, and a one-parameter group of symplectic automorphisms $\varphi_t$ dictating time evolution. Observables are functions on $N$, with their algebraic structure governed by the Poisson bracket $\{f,g\} = \omega(\xi_f, \xi_g)$, where $\xi_f$ is the Hamiltonian vector field of $f$ satisfying $i_{\xi_f}\omega = -df$. The framework is illustrated by deriving Hamilton's equations for particle motion in Euclidean space, $\dot{x}=y, \dot{y}=-V'(x)/m$, from the Hamiltonian $H = \frac{my^2}{2} + V(x)$ and symplectic form $\omega = m \, dy \wedge dx$. Generalizations to particle motion on Riemannian manifolds, where the phase space is the tangent bundle $TM$, are also considered. \\

\noindent We then transition to variational principles, initially examining geodesics in Euclidean space via the energy functional $E(x) = \int \frac{1}{2}|\dot{x}|^2 \, dt$, whose stationarity yields the geodesic equation $\ddot{x}=0$. This principle is reformulated with greater generality and local insight using the calculus of differential forms on the space of paths $\mathcal{F} \times \mathbb{R}$. The Lagrangian 1-form $L = L_0(x,\dot{x}) \, dt \in \Omega^{0,1}$ and the variational 1-form $\gamma \in \Omega^{1,0}$ (e.g., $\gamma = m \langle \delta x, \dot{x} \rangle$ for a free particle) are introduced. The Euler-Lagrange equations are then rewritten locally as $\delta L + d\gamma = 0$ on the manifold of solutions, where $\delta$ and $d$ are exterior derivatives on $\mathcal{F}$ and $\mathbb{R}$ respectively. This framework naturally describes general Lagrangian systems defined by a Lagrangian function $L_0(x,\dot{x})$. \\

\noindent The deep connection between Lagrangian and Hamiltonian formalisms then discussed. The canonical 2-form $\omega = \delta\gamma \in \Omega^{2,0}$, derived from the Lagrangian data, is shown to be $\delta$-closed, $d$-closed (time-independent), and non-degenerate, thus defining the symplectic structure on the space of solutions. The Hamiltonian function $H = \iota_\xi(L+\gamma)$ is constructed as the Noether conserved quantity corresponding to time-translation invariance (where $\xi$ generates time evolution), satisfying $\delta H = -\iota_\xi\omega$. Finally, a geometric interpretation of $\gamma(t)$ as a connection 1-form on the space of solutions is presented. The relation $\delta \int_{t_0}^{t_1} L = -(\gamma(t_1) - \gamma(t_0))$ reveals that the integrated Lagrangian acts as an isomorphism between these connections. 

\end{abstract}

\setcounter{tocdepth}{3}
\localtableofcontents

\textbf{References}: For foundational treatments of classical mechanics, including Newtonian, Lagrangian, and Hamiltonian formalisms, Arnold's text \cite{Arnold1989Mechanics} provides a mathematically sophisticated perspective, while Feynman's lectures \cite{FeynmanLeastAction} offer valuable physical intuition, particularly regarding the principle of least action. 

\newpage

\subsection{Hamiltonian Systems}
\label{sec:hamiltonian_systems}

We now introduce the central objects of study in Hamiltonian mechanics: Hamiltonian systems. These provide a powerful geometric framework for describing the evolution of many physical systems, from classical particle mechanics to more abstract dynamical systems. The main object of this framework is a symplectic manifold, whose structure dictates the rules of motion.

\subsubsection{Basic Definitions}
\label{subsec:hamiltonian_def_structure}

\begin{definition}[Hamiltonian System]
\,
\begin{enumerate}
    \item A \textbf{Hamiltonian system} consists of a symplectic manifold $(N, \omega)$ and a 1-parameter group $\varphi_{t}$ of symplectic automorphisms of $N$.\footnote{A symplectic automorphism is a diffeomorphism $\varphi: N \rightarrow N$ such that $\varphi^*\omega = \omega$.} Here, $N$ is the \textbf{phase space} and $\omega$ is the symplectic form.
    \item A Hamiltonian system is \textbf{free} if $(N,\omega)$ is an affine symplectic space and $\varphi_{t}$ is a 1-parameter group of affine symplectic automorphisms.
\end{enumerate}
\end{definition}

The symplectic manifold $N$ represents the space of possible states of the system, and the symplectic form $\omega$ gives the Poisson bracket structure of the system. The 1-parameter group of symplectic automorphisms $\varphi_{t}$ describes the time evolution of the system, mapping an initial state at time $t=0$ to its state at time $t$.

A Hamiltonian system naturally gives rise to a mechanical system:

\begin{enumerate}
    \item The space of pure states is the symplectic manifold $N$ itself. The space of all states can be considered as the space of probability distributions on $N$, denoted by $\text{Prob}(N)$.

    \item The algebra of observables is given by the space of complex-valued Borel functions on $N$, denoted by $\mathcal{O} = \text{Borel}(N; \C)$. The real structure is given by complex conjugation. The subspace of smooth functions $\mathcal{O}^{\infty} = C^{\infty}(N; \C)$ is dense in $\mathcal{O}$ and carries a Lie algebra structure defined by the Poisson bracket:
    \[ \{f, g\} = \omega(\xi_f, \xi_g) \]
    where $\xi_f$ and $\xi_g$ are the symplectic gradient vector fields of the functions $f$ and $g$, respectively, defined by $df = \omega(\xi_f, \cdot)$ and $dg = \omega(\xi_g, \cdot)$. The real observables are given by the real-valued Borel functions $\mathcal{O}_{\R} = \text{Borel}(N; \R)$, with a dense subspace of smooth real-valued functions $\mathcal{O}_{\R}^{\infty} = C^{\infty}(N; \R) = \Omega_N^0$.
    
    \item The measurement pairing associates an observable $A: N \rightarrow \R$ and a state $\sigma \in \text{Prob}(N)$ to the pushforward probability measure $A_*\sigma \in \text{Prob}(\R)$, which describes the probability distribution of the measured values of the observable $A$ in the state $\sigma$.
    
    \item A smooth observable $f \in \Omega_N^0$ has a symplectic gradient vector field $\xi_f \in \mathcal{X}_N$ defined by $df = \omega(\xi_f, \cdot)$. If $\xi_f$ is complete (i.e., its integral curves exist for all time), it generates a 1-parameter flow $\psi_t$ of symplectic diffeomorphisms. This flow induces a flow on states (by pushing forward the measure) and on complex observables (by pulling back the function), and these induced flows are compatible.
    
    \item The chosen 1-parameter group $\varphi_t$ of symplectic automorphisms in the definition of a Hamiltonian system is precisely the flow that generates the time evolution of the system on both states and observables.
\end{enumerate}

\begin{remark}
It is important to note that while every smooth observable has a symplectic gradient vector field, this vector field may not always generate a global flow. The existence of a global flow requires completeness of the vector field, which is not guaranteed in general. Therefore, the axiom regarding the existence of a global flow for all smooth observables is often something we hope for, rather than an axiom.
\end{remark}

\subsubsection{Particle Motion in Euclidean Space}

We consider the motion of a particle on the Euclidean line $\mathbb{E}^1$. The dynamics of such a system are described within the Hamiltonian formalism.

Let $N$ be the space of classical trajectories, meaning the space of functions $x: \mathbb{R} \to \mathbb{E}^1$ that satisfy Newton's second law for a given potential. For any specific trajectory, given an initial time $t_0 \in \mathbb{R}$ and initial conditions (position and velocity), a local solution for the trajectory exists. It is assumed that the potential $V$ is sufficiently well-behaved, ensuring that solutions $x(t)$ exist for all time.

Under these conditions, for any chosen time $t_0 \in \mathbb{R}$, the evaluation map
\[
N \longrightarrow \mathbb{E}^1 \times \mathbb{R}
\]
which maps a trajectory $x(t)$ to its position $x(t_0)$ and velocity $\dot{x}(t_0)$ at time $t_0$, is a diffeomorphism. This diffeomorphism allows an identification of the space of trajectories $N$ with the phase space $\mathbb{E}^1 \times \mathbb{R}$.

We use coordinates $(x, y)$ for this phase space, where `$x$' is the coordinate on $\mathbb{E}^1$ representing position, and `$y$' is the coordinate on $\mathbb{R}$ representing velocity (i.e., $y=\dot{x}$). The symplectic form $\omega$ on this phase space is given by:
\[
\omega = m \, \mathrm{d}y \wedge \mathrm{d}x
\]
where $m$ is the mass of the particle. This symplectic form corresponds to the canonical form $\mathrm{d}p \wedge \mathrm{d}q$ under the identification $q=x$ (position) and $p=my$ (canonical momentum).

The time evolution flow $\varphi_t: N \to N$ is generated by the Hamiltonian function $H$. For this system, $H$ is the total mechanical energy, expressed as the sum of kinetic and potential energy:
\[
H(x,y) = \frac{my^2}{2} + V(x)
\]
where $V(x)$ is the potential energy function. The exterior derivative (or differential) of $H$ is:
\[ \mathrm{d}H = \frac{\partial H}{\partial x} \mathrm{d}x + \frac{\partial H}{\partial y} \mathrm{d}y = V'(x) \mathrm{d}x + my \, \mathrm{d}y \]
where $V'(x) = \frac{\mathrm{d}V}{\mathrm{d}x}$.

The dynamics are governed by the Hamiltonian vector field $\xi_H$ associated with $H$. The symplectic gradient is defined by the relation:
\[ i_{\xi_H}\omega = -\mathrm{d}H \]
where $i_{\xi_H}\omega$ is the interior product of $\omega$ with $\xi_H$. Let the components of $\xi_H$ in our coordinate system be $A(x,y)$ and $B(x,y)$, such that:
\[ \xi_H = A(x,y) \frac{\partial}{\partial x} + B(x,y) \frac{\partial}{\partial y} \]
We compute the interior product $i_{\xi_H}\omega$:
\begin{align*}
i_{\xi_H}\omega &= i_{A \frac{\partial}{\partial x} + B \frac{\partial}{\partial y}} (m \, \mathrm{d}y \wedge \mathrm{d}x) \\
&= A \cdot i_{\frac{\partial}{\partial x}} (m \, \mathrm{d}y \wedge \mathrm{d}x) + B \cdot i_{\frac{\partial}{\partial y}} (m \, \mathrm{d}y \wedge \mathrm{d}x)
\end{align*}
To evaluate the interior products with the basis vector fields:
\begin{align*}
i_{\frac{\partial}{\partial x}} (m \, \mathrm{d}y \wedge \mathrm{d}x) &= m \left( (\mathrm{d}y (\frac{\partial}{\partial x})) \mathrm{d}x - (\mathrm{d}x (\frac{\partial}{\partial x})) \mathrm{d}y \right) \\
&= m \left( 0 \cdot \mathrm{d}x - 1 \cdot \mathrm{d}y \right) = -m \, \mathrm{d}y
\end{align*}
and
\begin{align*}
i_{\frac{\partial}{\partial y}} (m \, \mathrm{d}y \wedge \mathrm{d}x) &= m \left( (\mathrm{d}y (\frac{\partial}{\partial y})) \mathrm{d}x - (\mathrm{d}x (\frac{\partial}{\partial y})) \mathrm{d}y \right) \\
&= m \left( 1 \cdot \mathrm{d}x - 0 \cdot \mathrm{d}y \right) = m \, \mathrm{d}x
\end{align*}
Substituting these results back into the expression for $i_{\xi_H}\omega$:
\[ i_{\xi_H}\omega = A(-m \, \mathrm{d}y) + B(m \, \mathrm{d}x) = mB \, \mathrm{d}x - mA \, \mathrm{d}y \]
Now, we apply the defining equation $i_{\xi_H}\omega = -\mathrm{d}H$:
\[ mB \, \mathrm{d}x - mA \, \mathrm{d}y = -(V'(x) \mathrm{d}x + my \, \mathrm{d}y) \]
\[ mB \, \mathrm{d}x - mA \, \mathrm{d}y = -V'(x) \mathrm{d}x - my \, \mathrm{d}y \]
For this equality of 1-forms to hold, the coefficients of the basis 1-forms ($\mathrm{d}x$ and $\mathrm{d}y$) on both sides must be equal:
\begin{itemize}
    \item Equating coefficients of $\mathrm{d}x$: \quad $mB = -V'(x) \implies B = -\frac{1}{m}V'(x)$
    \item Equating coefficients of $\mathrm{d}y$: \quad $-mA = -my \implies A = y$
\end{itemize}
Thus, the Hamiltonian vector field is explicitly:
\[
\xi_H = y \frac{\partial}{\partial x} - \frac{1}{m}V'(x) \frac{\partial}{\partial y} 
\]
Consequently, Hamilton's equations for the flow $\varphi_t$, which describe the time evolution of $x$ and $y$, are obtained from the components of $\xi_H$ (i.e., $\dot{x} = A$ and $\dot{y} = B$):
\[
\left\{
\begin{aligned}
\dot{x} &= y \\
\dot{y} &= -\frac{1}{m}V'(x)
\end{aligned}
\right.
\]
These first-order differential equations are equivalent to Newton's second law, $m\ddot{x} = -V'(x)$. This can be verified by noting that $\dot{x}=y$, so differentiating with respect to time yields $\ddot{x}=\dot{y}$. Substituting this into the second Hamilton's equation gives $\ddot{x} = -\frac{1}{m}V'(x)$, or $m\ddot{x} = -V'(x)$.

\subsubsection{Remarks on Hamiltonian Systems}

\begin{remark}
\,
\begin{enumerate}
    \item The transition from Newton's second-order differential equation to the Hamiltonian system of two first-order differential equations is a standard technique achieved by introducing the velocity (or momentum) as an independent variable.
    \item The phase space $N = \mathbb{E}^1 \times \R$ is an affine space. For a free particle ($V=0$), the Hamiltonian is $H = \frac{1}{2}my^2$, and the equations of motion are $\dot{x} = y$, $\dot{y} = 0$, which represent a flow by affine symplectic diffeomorphisms. For a harmonic oscillator ($V(x) = \frac{1}{2}kx^2$ with $k > 0$), the Hamiltonian is $H = \frac{1}{2}my^2 + \frac{1}{2}kx^2$, and the equations of motion are $\dot{x} = y$, $\dot{y} = -\frac{k}{m}x$, which also generate a flow by affine symplectic diffeomorphisms. Thus, both systems are free Hamiltonian systems.
    \item Our example be generalized to the motion of a particle on a Riemannian manifold $(M, g)$ under the influence of a potential $V: M \rightarrow \R$. In this case, the acceleration $\ddot{x}$ is replaced by the acceleration computed using the Levi-Civita covariant derivative associated with the metric $g$. The phase space $N$ becomes the tangent bundle $TM$ of the configuration space $M$. Assuming that the equations of motion have solutions defined for all time, the map that takes a trajectory to its position and velocity at a given time is a diffeomorphism onto $TM$. The Riemannian metric $g$ induces an isomorphism between the tangent bundle $TM$ and the cotangent bundle $T^*M$. The symplectic form on $TM$ is related to the pullback of the tautological symplectic form on $T^*M$. Specifically, using the metric to identify $TM$ with $T^*M$, the symplectic form is $m$ times the pullback of the tautological 2-form. The Hamiltonian function is given by:
    \[ H(\xi) = \frac{m}{2} |\xi|^2 + V(\pi(\xi)) \]
    where $\xi \in TM$, $|\xi|^2 = g(\xi, \xi)$ is the squared norm of the tangent vector with respect to the metric $g$, and $\pi: TM \rightarrow M$ is the projection map. The Riemannian manifold $M$ is often referred to as the configuration space, while $N = TM$ is the phase space.
    \item Later in this lecture, we will demonstrate how to derive the symplectic form $\omega = m \, dy \wedge dx$ and the Hamiltonian function $H = \frac{1}{2}my^2 + V(x)$ (and their generalizations to Riemannian manifolds) from a more fundamental starting point: the Lagrangian formalism.
\end{enumerate}
\end{remark}

\begin{example}
The phase space $N$ of a Hamiltonian system does not necessarily have to be the tangent bundle of a manifold, and the 1-parameter group $\varphi_t$ of symplectic automorphisms need not always be generated by a global Hamiltonian function. Consider the 2-torus $N = \R / 2\pi\Z \times \R / 2\pi\Z$ with coordinates $(\theta^1, \theta^2)$. The standard symplectic form on $\R^2$ is $\omega_0 = dx \wedge dy$. On the torus, we can consider the symplectic form $\omega = d\theta^1 \wedge d\theta^2$. The 1-parameter group of diffeomorphisms given by:
\[ \varphi_t(\theta^1, \theta^2) = (\theta^1 + t, \theta^2) \]
is a symplectic automorphism because:
\[ \varphi_t^* (d\theta^1 \wedge d\theta^2) = d(\theta^1 + t) \wedge d\theta^2 = d\theta^1 \wedge d\theta^2 \]
This system can be seen as the quotient of a free system on $\R^2$ (with Hamiltonian $H(x, y) = y$) by the action of the discrete group $(2\pi\Z \times 2\pi\Z)$ of automorphisms, where the quotient map is $(x, y) \mapsto (x \mod 2\pi, y \mod 2\pi)$.
\end{example}

\subsection{Variational Principles for Geodesics}

\subsubsection{Introduction}

Variational principles play a fundamental role in many areas of mathematics and physics, providing an alternative way to formulate equations of motion and geometric structures. They are particularly important in global analysis on manifolds, where they are used to study the existence and properties of critical points of functionals defined on infinite-dimensional spaces of maps.

Let $V$ be a finite-dimensional real inner product space, and let $A$ be an affine space over $V$ (a Euclidean space). Fix real numbers $a < b$, and consider smooth paths $x: [a, b] \rightarrow A$. The length of the path is given by:
\[ L(x) = \int_a^b |\dot{x}(t)| \, dt = \int_a^b \sqrt{\langle \dot{x}(t), \dot{x}(t) \rangle} \, dt \]
where $\dot{x}(t)$ is the velocity vector of the path at time $t$, and $|\cdot|$ denotes the norm induced by the inner product $\langle \cdot, \cdot \rangle$. The length of the path is independent of the parametrization.

Instead of directly working with the length functional, it is often more convenient to consider the energy functional:
\[ E(x) = \int_a^b \frac{1}{2} |\dot{x}(t)|^2 \, dt = \int_a^b \frac{1}{2} \langle \dot{x}(t), \dot{x}(t) \rangle \, dt \]
A path that minimizes the energy functional (among paths with fixed endpoints) also minimizes the length functional, and vice versa, provided the length is non-zero.

To find the paths that make the energy functional stationary, we consider a variation of a given path $x(t)$. Let $x(t, s)$ be a smooth family of paths parametrized by $s \in (-\epsilon, \epsilon)$, such that $x(t, 0) = x(t)$ and the endpoints are fixed, i.e., $x(a, s) = x(a, 0)$ and $x(b, s) = x(b, 0)$ for all $s$. Let $\delta x(t) = \frac{\partial x}{\partial s}(t, 0)$ be the variation vector field along the path $x(t)$.

The variation of the energy functional is given by:
\[ \delta E = \frac{d}{ds} \Big|_{s=0} E(x(\cdot, s)) = \frac{d}{ds} \Big|_{s=0} \int_a^b \frac{1}{2} \langle \frac{\partial x}{\partial t}(t, s), \frac{\partial x}{\partial t}(t, s) \rangle \, dt \]
We can interchange the derivative with respect to $s$ and the integral:
\[ \delta E = \int_a^b \frac{1}{2} \frac{\partial}{\partial s} \langle \frac{\partial x}{\partial t}(t, s), \frac{\partial x}{\partial t}(t, s) \rangle \Big|_{s=0} \, dt \]
\[ = \int_a^b \langle \frac{\partial^2 x}{\partial s \partial t}(t, 0), \frac{\partial x}{\partial t}(t, 0) \rangle \, dt \]
Since the partial derivatives commute, $\frac{\partial^2 x}{\partial s \partial t} = \frac{\partial^2 x}{\partial t \partial s} = \frac{\partial}{\partial t} (\frac{\partial x}{\partial s}) = \frac{\partial}{\partial t} (\delta x) = \delta \dot{x}$. Thus,
\[ \delta E = \int_a^b \langle \delta \dot{x}(t), \dot{x}(t) \rangle \, dt \]
Using integration by parts, we have:
\[ \int_a^b \langle \delta \dot{x}(t), \dot{x}(t) \rangle \, dt = \Big[ \langle \delta x(t), \dot{x}(t) \rangle \Big]_a^b - \int_a^b \langle \delta x(t), \ddot{x}(t) \rangle \, dt \]
Since the endpoints are fixed, $\delta x(a) = 0$ and $\delta x(b) = 0$, so the boundary term vanishes. Therefore,
\[ \delta E = - \int_a^b \langle \delta x(t), \ddot{x}(t) \rangle \, dt \]
For the energy functional to be stationary ($\delta E = 0$) for all variations $\delta x$ with fixed endpoints, we must have $\ddot{x}(t) = 0$ for all $t \in [a, b]$. This equation describes a motion with constant velocity, which is a geodesic in the Euclidean space $A$.

\subsubsection{Via Differential Forms}

The variational principle, which yields the equations of motion such as the geodesic equation for a free particle, possesses an important local structure. This fundamental locality is most effectively revealed by re-studying the variational calculus in the language of differential forms on spaces of paths.

\begin{remark}[On Boundary Conditions and Locality]
\label{rem:boundary_locality}
In the computation of the first variation of the energy functional $E(x)$, one encounters a boundary term arising from integration by parts:
\begin{equation}
\int_{[a,b]} d_t \langle \delta x(t), \dot{x}(t) \rangle \, dt = \langle \delta x, \dot{x} \rangle|_{t=b} - \langle \delta x, \dot{x} \rangle|_{t=a}.
\end{equation}
Let us denote $\gamma_0(t) = \langle \delta x(t), \dot{x}(t) \rangle$. This $\gamma_0(t)$ represents, for each time $t$, a function on the space of variations $S$ (or, more generally, on the infinite-dimensional manifold of paths $\mathcal{F}$). The full variational 1-form $\gamma$ discussed later will be $\gamma = \gamma_0 \in \Omega^{1,0}$. It is customary in variational problems to impose boundary conditions on the paths $x(t)$ such that these boundary terms vanish. For instance, one might fix the initial and final positions, $x(a)$ and $x(b)$, implying $\delta x(a) = \delta x(b) = 0$.

Importantly, however, the Euler-Lagrange equation itself (in this case, the geodesic equation $\ddot{x}(t)=0$) is a local condition in time. Its derivation relies on the vanishing of the integrand of the variation for arbitrary variations $\delta x(t)$ supported away from the boundary. Thus, the specific choice of boundary conditions, while important for selecting a particular solution, does not affect the form of the local equations of motion.
\end{remark}

To study this local structure more clearly, we transition our focus from the global energy functional to the local \textit{Lagrangian function} $L_0(x,\dot{x}) = \frac{1}{2}\langle \dot{x}, \dot{x} \rangle$ and its associated \textit{Lagrangian 1-form} (a 1-form with respect to the time variable)
\begin{equation}
\label{eq:lagrangian_form}
L = L_0(x,\dot{x}) \, dt = \frac{1}{2}\langle \dot{x}, \dot{x} \rangle \, dt.
\end{equation}
Let $S$ be a smooth manifold parameterizing a family of paths $x: S \times [a,b] \to A$. Then $L$ can be viewed as an element of $\Omega^{0,1}(S \times [a,b])$, i.e., a 0-form on $S$ and a 1-form on $[a,b]$. The total energy is the pushforward $E = (pr_S)_*(L) \in \Omega^0(S)$, where $pr_S: S \times [a,b] \to S$ is the projection.

The natural arena for analyzing the local structure of variational problems is the product space $S \times [a,b]$ (or, more generally, $\mathcal{F} \times [a,b]$, where $\mathcal{F}$ denotes the infinite-dimensional manifold of smooth paths $x: [a,b] \to A$). The calculus of variations can be elegantly formulated using the de Rham complex on this product space. We distinguish between variations of paths and evolution in time by introducing two differential operators:
\begin{itemize}
    \item $\delta$: the exterior derivative acting on forms with respect to the variables parameterizing the space of paths $S$ (or $\mathcal{F}$).
    \item $d$: the exterior derivative acting on forms with respect to the time variable $t \in [a,b]$, so $d = dt \frac{d}{dt}$.
\end{itemize}
The total exterior derivative on $\Omega^\bullet(S \times [a,b])$ is $D = \delta + d$. These operators naturally endow the space of differential forms with a bicomplex structure, denoted $\Omega^{\bullet,\bullet}(S \times [a,b])$. A form $\omega \in \Omega^{p,q}(S \times [a,b])$ has degree $p$ in the $S$-variables (often called the \dq{variational degree}) and degree $q$ in the time variable.

Within this bicomplex framework:
\begin{itemize}
    \item The Lagrangian 1-form $L = L_0(x,\dot{x}) \, dt$ is manifestly an element of $\Omega^{0,1}(S \times [a,b])$.
    \item The variational 1-form $\gamma = \langle \delta x, \dot{x} \rangle$ (where $\delta x$ is the 1-form on $S$ obtained by applying $\delta$ to $x$) is an element of $\Omega^{1,0}(S \times [a,b])$. It is a 1-form with respect to variations on $S$ and a 0-form in time.
\end{itemize}
Applying the exterior derivative $\delta$ to $L$ yields $\delta L = (\delta L_0) \, dt = \langle \dot{x}, \delta \dot{x} \rangle \, dt \in \Omega^{1,1}(S \times [a,b])$. Similarly, applying $d$ to $\gamma$ gives $d\gamma = \left(\frac{d}{dt}\langle \delta x, \dot{x} \rangle\right) dt \in \Omega^{1,1}(S \times [a,b])$.
The important identity derived from the standard integration by parts procedure is:
\begin{equation} \label{eq:IBP_identity}
\langle \dot{x}, \delta \dot{x} \rangle = -\langle \delta x, \ddot{x} \rangle - \frac{d}{dt}\langle \delta x, \dot{x} \rangle.
\end{equation}
Multiplying by $dt$, this identity translates into the language of forms as:
$ \delta L = -\langle \delta x, \ddot{x} \rangle dt - d\gamma. $
This equation can be rewritten as $\delta L + d\gamma = -\langle \delta x, \ddot{x} \rangle dt$. The term $EL[\delta x] := -\langle \delta x, \ddot{x} \rangle dt$ is the Euler-Lagrange form associated with the variation $\delta x$. The requirement that $EL[\delta x]$ vanishes for arbitrary variations $\delta x$ (supported away from boundaries) yields the Euler-Lagrange equation $\ddot{x}=0$.

Let $N$ be the manifold of solutions to the equations of motion (i.e., paths $x(t)$ satisfying $\ddot{x}(t)=0$). These solutions are typically considered to be defined for all $t \in \mathbb{R}$. On the restricted space $N \times \mathbb{R}$, where $\ddot{x}=0$ by definition, the Euler-Lagrange form $EL[\delta x]$ vanishes. Consequently, the relationship $\delta L + d\gamma = EL[\delta x]$ simplifies to:
\begin{equation}
\label{eq:local_EL}
\delta L + d\gamma = 0 \quad \text{on } N \times \mathbb{R}.
\end{equation}
Here, $L \in \Omega^{0,1}(N \times \mathbb{R})$ and $\gamma \in \Omega^{1,0}(N \times \mathbb{R})$. Equation \eqref{eq:local_EL} expresses the vanishing of a $(1,1)$-form on the space of solutions extended over time. It is a very important piece of the covariant/geometric formulation of Lagrangian mechanics and field theory.

Furthermore, for the free particle Lagrangian, $L_0 = \frac{1}{2}|\dot{x}|^2$, $dL = d_t (\frac{1}{2}|\dot{x}|^2 dt) = \left( \frac{d}{dt}(\frac{1}{2}|\dot{x}|^2) \right) dt \wedge dt = 0$, since $dt \wedge dt = 0$. If one also assumes that $\gamma$ is $\delta$-closed on $N$ (i.e., $\delta\gamma = 0$), which is often the case for forms constructed from fields and their first variations without involving higher derivatives of variations, then the combined object $L+\gamma$ is $D$-closed on $N \times \mathbb{R}$:
$D(L+\gamma) = \delta L + dL + \delta\gamma + d\gamma = (\delta L + d\gamma) + dL + \delta\gamma = 0 + 0 + 0 = 0.$
The condition $\delta L + d\gamma = 0$ is thus the important local statement encapsulating the equations of motion within this formalism.

\begin{remark}[Further Aspects of the Variational Formalism]
\label{rem:further_aspects}
\,
\begin{enumerate}
    \item \textit{Origin of the variational 1-form $\gamma$:} The 1-form $\gamma = \langle \delta x, \dot{x} \rangle$ arises naturally from the integration by parts procedure. A more sophisticated geometric perspective, often developed in the context of the multisymplectic formalism or the calculus of variations on jet bundles, provides a construction of such forms (and their generalizations) without explicit reliance on integration. This involves defining appropriate currents and forms on spaces of field histories.

    \item \textit{A priori justification for the locality of Euler-Lagrange equations:} The fact that the critical point condition for an action functional $E(x) = \int_a^b L_0(x(t), \dot{x}(t)) dt$ yields a differential equation (local in time) can be understood intuitively. If a path $x_0$ minimizes (or extremizes) $E(x)$, it must necessarily do so for any sub-interval. Consider an infinitesimal interval $[t_0, t_0 + \varepsilon]$ for small $\varepsilon > 0$. The condition that $x_0$ extremizes $\int_{t_0}^{t_0+\varepsilon} L_0 dt$ for all such infinitesimal intervals, upon taking the limit $\varepsilon \to 0$ (after appropriate scaling), leads directly to the Euler-Lagrange equations evaluated at $t_0$. This underscores the local nature of these equations.

    \item \textit{Significance of variational principles:} The formulation of physical laws or geometric problems as variational principles is exceptionally powerful. If a system of ordinary or partial differential equations can be identified as the Euler-Lagrange equations for some functional (the action or energy), this functional provides important insights.
    \begin{itemize}
        \item It often serves as a guiding principle for constructing numerical methods to find solutions, e.g., via direct minimization or gradient flow techniques.
        \item The functional itself can yield important \textit{a priori} estimates (bounds, coercivity properties, etc.) essential for the analytical study of the existence, uniqueness, and regularity of solutions to the differential equations.
        \item Symmetries of the action functional, via Noether's theorem, lead to conservation laws, which are fundamental in physical theories.
    \end{itemize}
\end{enumerate}
\end{remark}


\subsection{Lagrangian Systems}

The Lagrangian approach to mechanics, rooted in the principle of stationary action (often attributed to Maupertuis, with significant contributions from Euler, Leibniz, and others), provides a powerful alternative to the Newtonian framework. While its historical development is rich, our focus here is on the formal mathematical structure, which underpins many areas of modern physics and geometry.

\begin{remark}
\label{rem:lagrangian_advantages}
\,
\begin{enumerate}
    \item \textit{Scope:} It is noteworthy that not every Hamiltonian system admits a Lagrangian formulation. The existence of a Lagrangian description often implies additional structure or regularity.
    \item \textit{Analytical Power:} A Lagrangian formulation can lead to powerful estimates and greater analytical control over the system's dynamics, particularly when studying the existence and properties of solutions to the equations of motion.
    \item \textit{Geometric Structure:} As we will explore, the Lagrangian viewpoint often reveals deeper geometric structures associated with a mechanical system, including connections to symplectic and contact geometry.
    \item \textit{Quantum Transition:} The Lagrangian, specifically the action functional, is the fundamental starting point for Feynman's path integral formulation of quantum mechanics, providing a direct bridge from classical to quantum descriptions.
\end{enumerate}
\end{remark}

\subsubsection{The Principle of Stationary Action and Equations of Motion}

Consider a particle of mass $m$ moving on a Riemannian manifold $(M,g)$, subject to a potential energy function $V: M \to \R$. A path of the particle is a smooth map $x: [a,b] \to M$. The \textit{Lagrangian} $L_0: TM \to \R$ for this system is defined as the difference between kinetic and potential energy:
\begin{equation}
    L_0(x(t), \dot{x}(t)) = \frac{m}{2} g(\dot{x}(t), \dot{x}(t)) - V(x(t)),
\end{equation}
where $\dot{x}(t) \in T_{x(t)}M$ is the velocity vector. The \textit{action functional} $S$ assigns to each path $x$ the integral of the Lagrangian over the time interval:
\begin{equation}
\label{eq:action_functional}
    S(x) = \int_a^b L_0(x(t), \dot{x}(t)) \, dt = \int_a^b \left[ \frac{m}{2} |\dot{x}(t)|_g^2 - V(x(t)) \right] dt.
\end{equation}
The principle of stationary action states that the actual paths taken by the system are those for which the action $S(x)$ is stationary with respect to variations of $x$ that fix the endpoints $x(a)$ and $x(b)$.

Both the Lagrangian $L_0$ and the action $S$ have physical dimensions of (mass)(length)$^2$(time)$^{-1}$. These are referred to as units of \textit{action}.

\begin{remark}[On the Nature of the Lagrangian]
\label{rem:nature_lagrangian}
The Lagrangian is a somewhat magical quantity. Unlike energy, force, or momentum, which often have direct and intuitive physical interpretations, the Lagrangian's role is more abstract, serving as the integrand whose stationary integral yields the equations of motion. Its remarkable ability in describing a vast range of physical phenomena makes it fundamentally important, even if its direct interpretation remains subtle.
\end{remark}

We employ the variational bicomplex formalism introduced previously. Let $\mathcal{F}$ be the infinite-dimensional manifold of smooth paths $x: \R \to M$. The relevant space for local analysis is $\mathcal{F} \times \R$.
The Lagrangian 1-form $L \in \Omega^{0,1}(\mathcal{F} \times \R)$ is $L = L_0(x,\dot{x}) dt$.
The variational 1-form $\gamma \in \Omega^{1,0}(\mathcal{F} \times \R)$ associated with variations $\delta x$ is given by the fiber derivative of $L_0$ with respect to velocity, paired with $\delta x$:
\begin{equation}
    \gamma = m \, g(\dot{x}, \delta x).
\end{equation}
The Euler-Lagrange equations, expressing the stationarity of the action, take the local form on the space $N \subset \mathcal{F}$ of solution paths:
\begin{equation}
\label{eq:EL_form_lagrangian}
    \delta L + d\gamma = 0 \quad \text{on } N \times \R.
\end{equation}
For the particle on $(M,g)$ with potential $V$, this abstract equation translates into Newton's second law, expressed in terms of the Levi-Civita covariant derivative $\nabla$ associated with $g$:
\begin{equation}
\label{eq:newtons_law_manifold}
    m \nabla_{\dot{x}}\dot{x} + \grad V(x) = 0,
\end{equation}
where $\grad V$ is the gradient vector field of $V$, defined by $g(\grad V, Y) = dV(Y)$ for any vector field $Y$ on $M$.

\subsubsection{Transition to Hamiltonian Mechanics}
The Lagrangian formalism provides a natural pathway to Hamiltonian mechanics, revealing the underlying symplectic structure of the phase space.

On the space of solutions $N \times \R$, we define the \textit{canonical 2-form} $\omega$ as the variation of $\gamma$:
\begin{definition}
The canonical 2-form $\omega$ is an element of $\Omega^{2,0}(N \times \R)$ given by
\begin{equation}
    \omega = \delta \gamma.
\end{equation}
\end{definition}
In the bicomplex notation, $\gamma$ is a $(1,0)$-form (a 1-form in variations, 0-form in time), and $\omega$ is a $(2,0)$-form (a 2-form in variations, 0-form in time). This means that for each fixed time $t$, $\omega_t = \omega|_{N \times \{t\}}$ is a 2-form on the space of solutions $N$.

\begin{lemma}[Properties of $\omega$]
\label{lemma:properties_omega}
The 2-form $\omega = \delta\gamma$ satisfies the following properties on $N \times \R$:
\begin{enumerate}
    \item $D\omega = 0$, which implies $\delta\omega = 0$ (i.e., $\omega$ is $\delta$-closed) and $d\omega = 0$ (i.e., $\omega$ is $d$-closed).
    \item The restriction $\omega_t = \omega|_{N \times \{t\}}$ is independent of time $t$.
    \item $\omega_t$ is nondegenerate on $N$ (viewed as a manifold of initial conditions).
\end{enumerate}
\end{lemma}
\begin{proof}
    \,
\begin{enumerate}
    \item Since $\omega = \delta\gamma$, we have $\delta\omega = \delta^2\gamma = 0$ because $\delta^2=0$.
    For $d\omega$, we use the identity $d\delta + \delta d = 0$ (which follows from $D^2 = (\delta+d)^2 = \delta^2+d^2+\delta d+d\delta = 0$) and the Euler-Lagrange equation $\delta L + d\gamma = 0$.
    $d\omega = d\delta\gamma = -\delta d\gamma$. From $\delta L + d\gamma = 0$, we have $d\gamma = -\delta L$.
    Thus, $d\omega = -\delta(-\delta L) = \delta^2 L = 0$.
    Since $D\omega = \delta\omega + d\omega$, we have $D\omega = 0+0=0$.

    \item The condition $d\omega = 0$ means that $\omega$ has no $dt$ component in its exterior derivative with respect to time. If we write $\omega = \sum_{i<j} \omega_{ij}(x,\delta x, t) \delta x^i \wedge \delta x^j$ locally (suppressing path indices for $N$), then $d\omega = \sum_{i<j} \frac{\partial \omega_{ij}}{\partial t} dt \wedge \delta x^i \wedge \delta x^j + \dots$. The vanishing of $d\omega$ implies $\frac{\partial \omega_{ij}}{\partial t} = 0$. Therefore, $\omega_t$ is independent of $t$.

    \item For the particle system, $\gamma = m \, g(\dot{x}, \delta x)$. Then
    \begin{equation}
        \omega = \delta \gamma = m \, \delta (g(\dot{x}, \delta x)) = m \, (g(\delta \dot{x}, \delta x) - g(\dot{x}, \delta^2 x)) = m \, g(\delta \dot{x}, \delta x)
    \end{equation}
    assuming a symmetric connection for $g$ such that $\delta$ and covariant differentiation commute appropriately, or more directly, $\omega = m ( \langle \delta\dot{x} \wedge \delta x \rangle_g )$ if interpreted on phase space. 
    
    More fundamentally, identifying the momentum $p_i = m g_{ij}\dot{x}^j$, $\gamma$ corresponds to $p_i \delta x^i$. Then $\omega = \delta p_i \wedge \delta x^i$, which is the canonical symplectic form on the cotangent bundle $T^*M$ (or a reduction thereof to $N$). This form is nondegenerate by definition of phase space.
\end{enumerate}
\end{proof}

\begin{corollary}
\label{cor:omega_symplectic}
For each $t \in \R$, the restriction $\omega_t = \omega|_{N \times \{t\}}$ is a symplectic form on $N$. Moreover, this symplectic structure is conserved in time.
\end{corollary}

The Hamiltonian function $H$ can be derived from the Lagrangian formalism, often interpreted as the conserved quantity associated with time-translation symmetry (Noether's theorem).

Let $\xi$ be the vector field on $\mathcal{F} \times \R$ that generates the negative of time translation. It is characterized by its action on $dt$ and variations $\delta x$: $\iota_{\xi}dt = -1$, and $\iota_{\xi}\delta x = \dot{x}$ when $\xi$ acts on $\gamma$.

The second condition means that the component of $\xi$ along $\mathcal{F}$ (the space of paths) corresponds to the velocity along the path when contracted with the variation 1-form $\delta x$. More precisely, if $\xi = \xi_{\mathcal{F}} - \partial_t$, then $\iota_{\xi_{\mathcal{F}}}\delta x = \dot{x}$.

The Hamiltonian $H$ is defined as a $(0,0)$-form (a function) on $\mathcal{F} \times \R$:
\begin{definition}
The Hamiltonian function $H$ is given by
\begin{equation}
\label{eq:def_hamiltonian}
    H = \iota_{\xi}(L+\gamma).
\end{equation}
\end{definition}

\begin{remark}[Noether's Theorem and the Hamiltonian]
\label{rem:noether_hamiltonian}
The definition of $H$ in \eqref{eq:def_hamiltonian} is directly motivated by Noether's theorem. If a Lagrangian system is invariant under a continuous symmetry, there is a corresponding conserved quantity. The Hamiltonian is precisely this conserved quantity associated with invariance under time translations. The term $L+\gamma$ can be seen as a component of a more general current.
\end{remark}

\begin{lemma}[Properties of the Hamiltonian]
\label{lemma:properties_hamiltonian}
On the space of solutions $N \times \R$:
\begin{enumerate}
    \item The Hamiltonian $H$ is conserved in time, i.e., $dH = 0$ (where $d$ is the time derivative component of $D$).
    \item The variation of $H$ is related to $\omega$ by $\delta H = -\iota_{\xi}\omega$.
\end{enumerate}
\end{lemma}
\begin{proof}
We use Cartan's magic formula $\mathcal{L}_{\xi} = D\iota_{\xi} + \iota_{\xi}D$.
Thus, $D H = D \iota_{\xi}(L+\gamma) = \mathcal{L}_{\xi}(L+\gamma) - \iota_{\xi}D(L+\gamma)$.
\begin{enumerate}
    \item For a time-invariant Lagrangian system, $L+\gamma$ is invariant under the flow of $\xi$ (time translation), so $\mathcal{L}_{\xi}(L+\gamma) = 0$.
    The Lagrangian 1-form $L=L_0 dt$ has $dL = d(L_0 dt) = (\frac{d L_0}{dt}) dt \wedge dt = 0$.
    On $N \times \R$, we have $\delta L + d\gamma = 0$.
    Therefore, $D(L+\gamma) = \delta L + dL + \delta\gamma + d\gamma = (\delta L+d\gamma) + dL + \delta\gamma = 0 + 0 + \omega = \omega$.
    So, $D H = 0 - \iota_{\xi}\omega = -\iota_{\xi}\omega$.
    Since $H$ is a function (a $(0,0)$-form), $DH = \delta H + dH$.
    The form $\omega = \delta\gamma$ is a $(2,0)$-form. Thus $\iota_{\xi}\omega$ is a $(1,0)$-form (it only has a $\delta$-component, no $dt$ component, as $\xi = \xi_{\mathcal{F}} - \partial_t$ and $\omega$ has no $dt$).
    Specifically, $\iota_{\xi}\omega = \iota_{\xi_{\mathcal{F}}}\omega - \iota_{\partial_t}\omega = \iota_{\xi_{\mathcal{F}}}\omega$ since $\omega$ is type $(2,0)$.
    Comparing degrees in $DH = -\iota_{\xi}\omega$:
    The $(0,1)$-component gives $dH = 0$, so $H$ is constant in time.

    \item The $(1,0)$-component gives $\delta H = -\iota_{\xi}\omega$. This is Hamilton's equation in form notation, relating the Hamiltonian to the symplectic form and the generator of time evolution.
\end{enumerate}
\end{proof}

\begin{example}[Particle on $\mathbb{E}^1$]
\label{ex:particle_E1_hamiltonian}
Consider a particle of mass $m$ moving on the Euclidean line $\mathbb{E}^1$ (so $g(\dot{x},\dot{x}) = \dot{x}^2$) with potential $V(x)$.
The Lagrangian 1-form is $L = \left(\frac{m}{2}\dot{x}^2 - V(x)\right)dt$.
The variational 1-form is $\gamma = m\dot{x}\delta x$.
So, $L+\gamma = \left(\frac{m}{2}\dot{x}^2 - V(x)\right)dt + m\dot{x}\delta x$.
Using $\iota_{\xi}dt = -1$ and $\iota_{\xi}\delta x = \dot{x}$ (meaning $\iota_{\xi_{\mathcal{F}}}\delta x = \dot{x}$):
\begin{align*}
    H &= \iota_{\xi}(L+\gamma) \\
      &= \iota_{\xi}\left(\left(\frac{m}{2}\dot{x}^2 - V(x)\right)dt\right) + \iota_{\xi}(m\dot{x}\delta x) \\
      &= -\left(\frac{m}{2}\dot{x}^2 - V(x)\right) + m\dot{x}(\iota_{\xi}\delta x) \\
      &= -\frac{m}{2}\dot{x}^2 + V(x) + m\dot{x}^2 \\
      &= \frac{m}{2}\dot{x}^2 + V(x).
\end{align*}
This is the familiar expression for the total energy (Hamiltonian) of the particle. Note the sign change of the potential $V(x)$ from the Lagrangian to the Hamiltonian.
\end{example}

\subsection{Geometric Interpretation of $\gamma$}

The variational approach to mechanics has deep connections to differential geometry and provides a powerful framework for understanding the underlying structures of physical theories.

On the space $N$ of classical trajectories, the equation $\delta L = -d\gamma$ holds, which implies $\delta \int_{t_0}^{t_1} L = -(\gamma(t_1) - \gamma(t_0))$. This equation has a deep geometric interpretation in terms of connections on principal bundles.

Consider $\R$ as a Lie group under addition. A connection on a trivial principal $\R$-bundle over $N$ is a real-valued 1-form on $N$. For each $t \in \R$, the 1-form $\gamma(t)$ can be interpreted as an $\R$-connection on $N \times \{t\}$. The equation $\delta \int_{t_0}^{t_1} L = -(\gamma(t_1) - \gamma(t_0))$ implies that addition by $\int_{t_0}^{t_1} L$ provides an isomorphism of trivial principal $\R$-bundles with connection, mapping the connection $\gamma(t_0)$ to $\gamma(t_1)$. These isomorphisms are compatible for different time intervals.

By taking invariant sections in the $\R$-direction, we obtain a principal $\R$-bundle $T \rightarrow N$ equipped with a connection whose curvature is the symplectic form $\omega = \delta \gamma$. This provides a new way to construct the symplectic form on $N$ as the curvature of a connection.

\begin{remark}
The first step in the path integral formulation of quantum theory involves exponentiating this principal $\R$-bundle with its connection. Quantization can be intuitively understood as an exponentiation process. The units of $\gamma$ are the units of action ($ML^2/T$). To perform the exponentiation, we need a constant with units of action, which is provided by Planck's constant $\hbar$. Applying the function $\exp(\frac{i}{\hbar} \cdot): \R \rightarrow \C^\times$, we can form a principal $\C^\times$-bundle with connection associated to the $\R$-bundle. The associated complex line bundle with covariant derivative is often called the prequantum line bundle. This construction is canonical from a Lagrangian formulation and highlights the additional geometric structure inherent in the Lagrangian formalism.
\end{remark}

\begin{example}[Particle on a Ring]
Consider a particle moving on a ring, which can be modeled as the Euclidean circle $\mathbb{E}^1 / 2\pi L \Z$ of length $2\pi L$. Let $m > 0$ be the mass of the particle. We define a 1-parameter family of Lagrangian theories parametrized by $\theta \in \R$:
\[ L = \left( \frac{m}{2} \dot{x}^2 + \frac{\theta}{2\pi} \dot{x} \right) dt \]
The first term is the usual kinetic energy, and the second term is locally exact but not globally exact on the circle since there is no global coordinate $x$ that covers the entire circle smoothly.

The space $N$ of classical solutions consists of motions with constant velocity. The Euler-Lagrange equation is $m \ddot{x} = 0$, so $x(t) = vt + x_0 \mod 2\pi L$. The phase space can be identified with the initial position and velocity, which are $(\theta^1, \theta^2) \in \R / 2\pi L \Z \times \R$. The symplectic form can be derived from the Lagrangian.

The Hamiltonian derived from this Lagrangian is independent of $\theta$. 
The Euler-Lagrange equation is $\frac{d}{dt} (m \dot{x} + \frac{\theta}{2\pi}) - 0 = 0$, so $m \ddot{x} = 0$. The solutions are $x(t) = vt + x_0 \mod 2\pi L$. The velocity $v$ is constant. The momentum is $p = m v + \frac{\theta}{2\pi}$. The Hamiltonian is $H = p \dot{x} - L = (m v + \frac{\theta}{2\pi}) v - (\frac{m}{2} v^2 + \frac{\theta}{2\pi} v) = \frac{m}{2} v^2 = \frac{1}{2m} (p - \frac{\theta}{2\pi})^2$.

The principal bundle with connection constructed from this Lagrangian depends on $\theta$. The connection 1-form is related to $\gamma = (m \dot{x} + \frac{\theta}{2\pi}) \delta x$. The curvature of this connection, which is the symplectic form, should be independent of $\theta$.

The dependence on $\theta$ manifests in the corresponding quantum systems, as will be seen in later discussions of quantization.
\end{example}

\begin{remark}
This example illustrates that different Lagrangian systems can give rise to the same Hamiltonian system. The Hamiltonian $H = \frac{p^2}{2m}$ on the phase space $T^* (\R / 2\pi L \Z) \cong (\R / 2\pi L \Z) \times \R$ (where $p$ is the momentum conjugate to the angular coordinate) corresponds to free motion on the ring, independent of $\theta$. However, the Lagrangian includes a term that affects the canonical momentum and the prequantum line bundle, leading to different quantum theories for different values of $\theta$. This highlights that Lagrangian systems can contain more information than Hamiltonian systems.
\end{remark}


